\newpage
\section{Método da seção crítica}

% Item 4 do roteiro: CLmax de asa limpa em baixa velocidade (mn 0.2) para
% cada CG, com e sem trimagem. Varrer alpha ate o cl_norm de alguma secao
% da asa cruzar o clmax de perfil. Salvar fs.txt de cada caso em avl/saidas/
% e gerar os graficos cl x y com as curvas limite de clmax.
% PREENCHER apos as rodadas.

O $C_{L\mathrm{max}}$ de asa limpa em baixa velocidade foi determinado pelo
método da seção crítica em $M = 0{,}2$. Para cada posição de CG, com e sem a
restrição de trimagem, o ângulo de ataque foi aumentado até que a
distribuição de $c_{\ell,\mathrm{norm}}$ de alguma estação da asa alcançasse
o limite de $c_{\ell\mathrm{max}}$ dos perfis obtido no Lab 03.

% Graficos cl x y (um por caso ou um combinado), com as curvas de limite:
% \begin{figure}[H] ... images/clxy_fwd_semtrim.png ... \end{figure}

\begin{table}[H]
\centering
\caption{Resultados do método da seção crítica para os quatro casos.}\label{tab:secao_critica}
\renewcommand{\arraystretch}{1.25}
\begin{tabular}{lccc}
\toprule
Caso & $\alpha_{\mathrm{max}}$ [graus] & $C_{L\mathrm{max}}$ & $\delta_e$ [graus] \\
\midrule
CG dianteiro sem trimagem & --- & --- & --- \\
CG dianteiro com trimagem & --- & --- & --- \\
CG traseiro sem trimagem & --- & --- & --- \\
CG traseiro com trimagem & --- & --- & --- \\
\bottomrule
\end{tabular}
\end{table}

% Discutir: onde comeca o estol (raiz ou ponta), se ha margem de aileron,
% e se as caracteristicas de estol sao adequadas.
